\section{Discussion}
\label{sec:discussion}

\subsection{Artistic Implications}

The depth-to-radius mapping (Eq.~\ref{eq:radius_mapping}) introduces a
new dimension of artistic control that goes beyond traditional
luminance-based stippling. By modulating dot size with scene depth, the
pipeline creates a tactile sense of three-dimensional structure within
the inherently flat polka dot medium---an effect that resonates with
Kusama's own exploration of infinity nets and immersive environments
where dots dissolve spatial boundaries. The Gaussian noise term
$\varepsilon$ is not merely a technical addition but an artistic
necessity: Kusama's dots, despite their geometric precision, are never
mechanically uniform. The controlled randomness replicates the
hand-painted quality that distinguishes her work from machine-generated
patterns.

The semantic-aware dual-color rendering introduces a computational
analogue to Kusama's practice of assigning different colored dots to
different surfaces within a single composition. In her installations,
walls, floors, and objects each receive dots in distinct colors,
creating a figure--ground dialogue. Our pipeline achieves a similar
effect by leveraging semantic segmentation to distinguish foreground
objects from background regions, assigning each a different dot color
within a shared palette.

\subsection{Limitations}

\paragraph{Semantic Taxonomy Mismatch.}
The YOLO26x-sem model is trained on the Cityscapes dataset, which was
designed for autonomous driving scenarios. Its 19-class taxonomy does
not include common artistic subjects such as animals, food, or
flowers. For the cat image, only 7.1\% of pixels were classified as
foreground (a small person region), and the pumpkin was classified as
``car'' rather than as a botanical object. While the dual-color
mechanism still functions---any non-background class triggers the
foreground color---the semantic distinction may not align with the
artistic intent. A model fine-tuned on a broader taxonomy (e.g.,
ADE20K \citep{zhou2017scene}) or on art-specific segmentation tasks
would improve semantic-aware rendering for general photographic
content.

\paragraph{Depth Estimation Accuracy.}
Monocular depth estimation is an inherently ill-posed problem, and the
estimated depth maps may contain errors, particularly for objects not
well-represented in the training data. These errors propagate to the
dot size assignment, potentially producing depth-inconsistent dot
patterns. However, since the pipeline uses relative depth ordering
rather than absolute metric values (the percentile clipping and
normalization in Eq.~\ref{eq:depth_norm} make the mapping robust to
scale errors), the visual impact of moderate depth estimation errors
is limited.

\paragraph{Computational Cost.}
The Voronoi stippling stage dominates the runtime at approximately
60\,s per image, primarily due to the iterative Lloyd relaxation. This
is acceptable for offline art generation but precludes real-time
interactive use. GPU-accelerated Voronoi computation
\citep{ray2018codesign} or approximate relaxation methods could reduce
this to interactive speeds.

\paragraph{Single-Image Input.}
The pipeline operates on single RGB images and does not exploit
multi-view or video input. For video sequences, temporal consistency
of dot placement and size would require additional mechanisms such as
optical flow-guided dot tracking or per-frame depth consistency
constraints.

\subsection{Future Directions}

We highlight two promising extensions.
First, the output of our pipeline could serve as a conditioning
image for a depth-aware ControlNet \citep{zhang2023adding} combined
with a Kusama-style LoRA \citep{hu2022lora}, producing a hybrid
result that combines algorithmic structural precision with learned
textural richness. Second, extending to video sequences via
optical flow-guided dot tracking and per-frame depth consistency
would enable temporally stable polka dot animations.
